% --- Archon physics formalization source begin ---
% archon:physics
% archon:covers PhyXMiniProblems/problem_phyx_mini_0234.lean
% archon:source-report reports/phyx_mini/problem_phyx_mini_0234.source.json

\chapter{Physics problem phyx\_mini\_0234}
\label{ch:PhyXMiniProblems_problem_phyx_mini_0234}

\paragraph{Problem source.}
\#\# Physical scenario

A light balloon filled with helium of density \textbackslash{}( 0.179\textbackslash{},\textbackslash{}mathrm\{kg/m\textasciicircum{}3\} \textbackslash{}) is tied to a light string of length \textbackslash{}( L = 3.00\textbackslash{},\textbackslash{}mathrm\{m\} \textbackslash{}). The string is tied to the ground forming an inverted simple pendulum (figure).The balloon is displaced slightly from equilibrium as in figure and released.Take the density of air to be \textbackslash{}( 1.20\textbackslash{},\textbackslash{}mathrm\{kg/m\textasciicircum{}3\} \textbackslash{}).Assume the air applies a buoyant force on the balloon but does not otherwise affect its motion.

\#\# Question

Determine the period of the motion.

\#\# Answer choices

- A: A: 1.76s

- B: B: 1.66s

- C: C: 1.56s

- D: D: 1.46s

\#\# Figure caption (auxiliary; use the image as primary evidence)

The image depicts a physical system involving two masses and three strings. Two vertical walls support the setup. The strings are arranged to form two triangles and a horizontal line. 

Objects and components:
- Two walls positioned on the left and right sides.
- Two spherical masses represented as blue balls, one above the other in the center.
- Three strings: two angled strings forming triangles with the horizontal string.

Relationships and positioning:
- The upper string runs horizontally between the walls with each half labeled "L."
- The two angled strings connect from the walls to meet at the upper mass, forming two right triangles with the horizontal string.
- The distance "y" is indicated vertically between the two masses.

Text labels:
- "L" denotes the lengths of the segments of the horizontal string.
- "y" represents the vertical distance between the two masses.

\#\# Dataset metadata

Wave Properties; Physical Model Grounding Reasoning, Multi-Formula Reasoning

\paragraph{Recorded answer/context.}
D: D: 1.46s

\paragraph{Figure/image path.}
/root/proposal\_for\_physic/science-mango/phyx\_mini\_run/phyx\_data/test\_image/234.png

\paragraph{Formalization target.}
create a compiling Lean file with sorry bodies at `PhyXMiniProblems/problem\_phyx\_mini\_0234.lean`. The Lean declarations must preserve the physical quantities, dimensions or dimensional roles, figure labels, governing-law hypotheses, and final relation expressed by this problem.
Use Mathlib/Physlib names found through LeanExplore where available. If a physics API is missing, introduce faithful local abstractions rather than scalar placeholder aliases.

\begin{theorem}[Physics formalization target]
\label{thm:physics:phyx_mini_0234:target}
\lean{PhyXMiniProblems.ProblemPhyXMini0234.balloonMotionPeriod_is_answer_D}
\uses{def:physics:phyx-mini-0234:phyxminiproblems-problemphyxmini0234-timeinseconds, def:physics:phyx-mini-0234:phyxminiproblems-problemphyxmini0234-balloonpendulumsetup, def:physics:phyx-mini-0234:phyxminiproblems-problemphyxmini0234-matchesverbalballoonscenario, def:physics:phyx-mini-0234:phyxminiproblems-problemphyxmini0234-matchessuppliedfigure, def:physics:phyx-mini-0234:phyxminiproblems-problemphyxmini0234-matchesproblemdata, def:physics:phyx-mini-0234:phyxminiproblems-problemphyxmini0234-usesstandardgravity, def:physics:phyx-mini-0234:phyxminiproblems-problemphyxmini0234-hasphysicalballoonparameters, def:physics:phyx-mini-0234:phyxminiproblems-problemphyxmini0234-satisfiesballoonpendulumlaws, def:physics:phyx-mini-0234:phyxminiproblems-problemphyxmini0234-satisfiesfirstordersmalloscillationmodel, def:physics:phyx-mini-0234:phyxminiproblems-problemphyxmini0234-answerchoiceinseconds, def:physics:phyx-mini-0234:phyxminiproblems-problemphyxmini0234-roundstonearesthundredth}
The assigned autoformalize agent should translate this physics problem into Lean declarations in the covered file, with theorem and lemma proof bodies written as `by sorry`.
\end{theorem}
\begin{proof}
This is an autoformalization task, not a proof task. Produce statements that can later be proved by the physics prover without weakening the physical model.
\end{proof}

% --- Archon declaration topology begin ---
\section{Declaration topology}
\begin{definition}[Length Quantity]
  \label{def:physics:phyx-mini-0234:phyxminiproblems-problemphyxmini0234-lengthquantity}
  \lean{PhyXMiniProblems.ProblemPhyXMini0234.LengthQuantity}
  A nonnegative physical length
\end{definition}
\begin{proof}
  This is the declared model definition of Length Quantity.
\end{proof}
\begin{definition}[Time Quantity]
  \label{def:physics:phyx-mini-0234:phyxminiproblems-problemphyxmini0234-timequantity}
  \lean{PhyXMiniProblems.ProblemPhyXMini0234.TimeQuantity}
  A nonnegative physical time interval
\end{definition}
\begin{proof}
  This is the declared model definition of Time Quantity.
\end{proof}
\begin{definition}[Volume Quantity]
  \label{def:physics:phyx-mini-0234:phyxminiproblems-problemphyxmini0234-volumequantity}
  \lean{PhyXMiniProblems.ProblemPhyXMini0234.VolumeQuantity}
  A nonnegative physical volume, with dimension length\textasciicircum{}3
\end{definition}
\begin{proof}
  This is the declared model definition of Volume Quantity.
\end{proof}
\begin{definition}[Mass Quantity]
  \label{def:physics:phyx-mini-0234:phyxminiproblems-problemphyxmini0234-massquantity}
  \lean{PhyXMiniProblems.ProblemPhyXMini0234.MassQuantity}
  A nonnegative physical mass
\end{definition}
\begin{proof}
  This is the declared model definition of Mass Quantity.
\end{proof}
\begin{definition}[Mass Density Quantity]
  \label{def:physics:phyx-mini-0234:phyxminiproblems-problemphyxmini0234-massdensityquantity}
  \lean{PhyXMiniProblems.ProblemPhyXMini0234.MassDensityQuantity}
  A nonnegative mass density, with dimension mass / length\textasciicircum{}3
\end{definition}
\begin{proof}
  This is the declared model definition of Mass Density Quantity.
\end{proof}
\begin{definition}[Acceleration Quantity]
  \label{def:physics:phyx-mini-0234:phyxminiproblems-problemphyxmini0234-accelerationquantity}
  \lean{PhyXMiniProblems.ProblemPhyXMini0234.AccelerationQuantity}
  A nonnegative acceleration magnitude
\end{definition}
\begin{proof}
  This is the declared model definition of Acceleration Quantity.
\end{proof}
\begin{definition}[Force Quantity]
  \label{def:physics:phyx-mini-0234:phyxminiproblems-problemphyxmini0234-forcequantity}
  \lean{PhyXMiniProblems.ProblemPhyXMini0234.ForceQuantity}
  A nonnegative force magnitude
\end{definition}
\begin{proof}
  This is the declared model definition of Force Quantity.
\end{proof}
\begin{definition}[Angular Frequency Quantity]
  \label{def:physics:phyx-mini-0234:phyxminiproblems-problemphyxmini0234-angularfrequencyquantity}
  \lean{PhyXMiniProblems.ProblemPhyXMini0234.AngularFrequencyQuantity}
  An angular frequency; radians are dimensionless
\end{definition}
\begin{proof}
  This is the declared model definition of Angular Frequency Quantity.
\end{proof}
\begin{definition}[Angular Velocity Quantity]
  \label{def:physics:phyx-mini-0234:phyxminiproblems-problemphyxmini0234-angularvelocityquantity}
  \lean{PhyXMiniProblems.ProblemPhyXMini0234.AngularVelocityQuantity}
  A signed angular velocity
\end{definition}
\begin{proof}
  This is the declared model definition of Angular Velocity Quantity.
\end{proof}
\begin{definition}[Angular Acceleration Quantity]
  \label{def:physics:phyx-mini-0234:phyxminiproblems-problemphyxmini0234-angularaccelerationquantity}
  \lean{PhyXMiniProblems.ProblemPhyXMini0234.AngularAccelerationQuantity}
  A signed angular acceleration
\end{definition}
\begin{proof}
  This is the declared model definition of Angular Acceleration Quantity.
\end{proof}
\begin{definition}[nonnegative SIReadout]
  \label{def:physics:phyx-mini-0234:phyxminiproblems-problemphyxmini0234-nonnegativesireadout}
  \lean{PhyXMiniProblems.ProblemPhyXMini0234.nonnegativeSIReadout}
  Read a nonnegative dimensionful quantity in coherent SI units
\end{definition}
\begin{proof}
  This is the declared model definition of nonnegative SIReadout.
\end{proof}
\begin{definition}[signed SIReadout]
  \label{def:physics:phyx-mini-0234:phyxminiproblems-problemphyxmini0234-signedsireadout}
  \lean{PhyXMiniProblems.ProblemPhyXMini0234.signedSIReadout}
  Read a signed dimensionful quantity in coherent SI units
\end{definition}
\begin{proof}
  This is the declared model definition of signed SIReadout.
\end{proof}
\begin{definition}[length In Meters]
  \label{def:physics:phyx-mini-0234:phyxminiproblems-problemphyxmini0234-lengthinmeters}
  \lean{PhyXMiniProblems.ProblemPhyXMini0234.lengthInMeters}
  \uses{def:physics:phyx-mini-0234:phyxminiproblems-problemphyxmini0234-lengthquantity, def:physics:phyx-mini-0234:phyxminiproblems-problemphyxmini0234-nonnegativesireadout}
  Metre readout of a physical length
\end{definition}
\begin{proof}
  This is the declared model definition of length In Meters.
\end{proof}
\begin{definition}[time In Seconds]
  \label{def:physics:phyx-mini-0234:phyxminiproblems-problemphyxmini0234-timeinseconds}
  \lean{PhyXMiniProblems.ProblemPhyXMini0234.timeInSeconds}
  \uses{def:physics:phyx-mini-0234:phyxminiproblems-problemphyxmini0234-timequantity, def:physics:phyx-mini-0234:phyxminiproblems-problemphyxmini0234-nonnegativesireadout}
  Second readout of a physical time interval
\end{definition}
\begin{proof}
  This is the declared model definition of time In Seconds.
\end{proof}
\begin{definition}[volume In Cubic Meters]
  \label{def:physics:phyx-mini-0234:phyxminiproblems-problemphyxmini0234-volumeincubicmeters}
  \lean{PhyXMiniProblems.ProblemPhyXMini0234.volumeInCubicMeters}
  \uses{def:physics:phyx-mini-0234:phyxminiproblems-problemphyxmini0234-volumequantity, def:physics:phyx-mini-0234:phyxminiproblems-problemphyxmini0234-nonnegativesireadout}
  Cubic-metre readout of a physical volume
\end{definition}
\begin{proof}
  This is the declared model definition of volume In Cubic Meters.
\end{proof}
\begin{definition}[mass In Kilograms]
  \label{def:physics:phyx-mini-0234:phyxminiproblems-problemphyxmini0234-massinkilograms}
  \lean{PhyXMiniProblems.ProblemPhyXMini0234.massInKilograms}
  \uses{def:physics:phyx-mini-0234:phyxminiproblems-problemphyxmini0234-massquantity, def:physics:phyx-mini-0234:phyxminiproblems-problemphyxmini0234-nonnegativesireadout}
  Kilogram readout of a physical mass
\end{definition}
\begin{proof}
  This is the declared model definition of mass In Kilograms.
\end{proof}
\begin{definition}[density In Kilograms Per Cubic Meter]
  \label{def:physics:phyx-mini-0234:phyxminiproblems-problemphyxmini0234-densityinkilogramspercubicmeter}
  \lean{PhyXMiniProblems.ProblemPhyXMini0234.densityInKilogramsPerCubicMeter}
  \uses{def:physics:phyx-mini-0234:phyxminiproblems-problemphyxmini0234-massdensityquantity, def:physics:phyx-mini-0234:phyxminiproblems-problemphyxmini0234-nonnegativesireadout}
  Kilogram-per-cubic-metre readout of a mass density
\end{definition}
\begin{proof}
  This is the declared model definition of density In Kilograms Per Cubic Meter.
\end{proof}
\begin{definition}[acceleration In Meters Per Second Squared]
  \label{def:physics:phyx-mini-0234:phyxminiproblems-problemphyxmini0234-accelerationinmeterspersecondsquared}
  \lean{PhyXMiniProblems.ProblemPhyXMini0234.accelerationInMetersPerSecondSquared}
  \uses{def:physics:phyx-mini-0234:phyxminiproblems-problemphyxmini0234-accelerationquantity, def:physics:phyx-mini-0234:phyxminiproblems-problemphyxmini0234-nonnegativesireadout}
  Metre-per-second-squared readout of an acceleration magnitude
\end{definition}
\begin{proof}
  This is the declared model definition of acceleration In Meters Per Second Squared.
\end{proof}
\begin{definition}[force In Newtons]
  \label{def:physics:phyx-mini-0234:phyxminiproblems-problemphyxmini0234-forceinnewtons}
  \lean{PhyXMiniProblems.ProblemPhyXMini0234.forceInNewtons}
  \uses{def:physics:phyx-mini-0234:phyxminiproblems-problemphyxmini0234-forcequantity, def:physics:phyx-mini-0234:phyxminiproblems-problemphyxmini0234-nonnegativesireadout}
  Newton readout of a force magnitude
\end{definition}
\begin{proof}
  This is the declared model definition of force In Newtons.
\end{proof}
\begin{definition}[angular Frequency In Radians Per Second]
  \label{def:physics:phyx-mini-0234:phyxminiproblems-problemphyxmini0234-angularfrequencyinradianspersecond}
  \lean{PhyXMiniProblems.ProblemPhyXMini0234.angularFrequencyInRadiansPerSecond}
  \uses{def:physics:phyx-mini-0234:phyxminiproblems-problemphyxmini0234-angularfrequencyquantity, def:physics:phyx-mini-0234:phyxminiproblems-problemphyxmini0234-nonnegativesireadout}
  Radian-per-second readout of an angular frequency
\end{definition}
\begin{proof}
  This is the declared model definition of angular Frequency In Radians Per Second.
\end{proof}
\begin{definition}[angular Velocity In Radians Per Second]
  \label{def:physics:phyx-mini-0234:phyxminiproblems-problemphyxmini0234-angularvelocityinradianspersecond}
  \lean{PhyXMiniProblems.ProblemPhyXMini0234.angularVelocityInRadiansPerSecond}
  \uses{def:physics:phyx-mini-0234:phyxminiproblems-problemphyxmini0234-angularvelocityquantity, def:physics:phyx-mini-0234:phyxminiproblems-problemphyxmini0234-signedsireadout}
  Radian-per-second readout of a signed angular velocity
\end{definition}
\begin{proof}
  This is the declared model definition of angular Velocity In Radians Per Second.
\end{proof}
\begin{definition}[angular Acceleration In Radians Per Second Squared]
  \label{def:physics:phyx-mini-0234:phyxminiproblems-problemphyxmini0234-angularaccelerationinradianspersecondsquared}
  \lean{PhyXMiniProblems.ProblemPhyXMini0234.angularAccelerationInRadiansPerSecondSquared}
  \uses{def:physics:phyx-mini-0234:phyxminiproblems-problemphyxmini0234-angularaccelerationquantity, def:physics:phyx-mini-0234:phyxminiproblems-problemphyxmini0234-signedsireadout}
  Radian-per-second-squared readout of a signed angular acceleration
\end{definition}
\begin{proof}
  This is the declared model definition of angular Acceleration In Radians Per Second Squared.
\end{proof}
\begin{definition}[Balloon Gas]
  \label{def:physics:phyx-mini-0234:phyxminiproblems-problemphyxmini0234-balloongas}
  \lean{PhyXMiniProblems.ProblemPhyXMini0234.BalloonGas}
  Gas filling the light balloon
\end{definition}
\begin{proof}
  This is the declared model definition of Balloon Gas.
\end{proof}
\begin{definition}[Tether Anchor]
  \label{def:physics:phyx-mini-0234:phyxminiproblems-problemphyxmini0234-tetheranchor}
  \lean{PhyXMiniProblems.ProblemPhyXMini0234.TetherAnchor}
  Where the end of the tether opposite the balloon is fixed
\end{definition}
\begin{proof}
  This is the declared model definition of Tether Anchor.
\end{proof}
\begin{definition}[Pendulum Orientation]
  \label{def:physics:phyx-mini-0234:phyxminiproblems-problemphyxmini0234-pendulumorientation}
  \lean{PhyXMiniProblems.ProblemPhyXMini0234.PendulumOrientation}
  Equilibrium orientation of the tethered balloon
\end{definition}
\begin{proof}
  This is the declared model definition of Pendulum Orientation.
\end{proof}
\begin{definition}[Auxiliary Mass Model]
  \label{def:physics:phyx-mini-0234:phyxminiproblems-problemphyxmini0234-auxiliarymassmodel}
  \lean{PhyXMiniProblems.ProblemPhyXMini0234.AuxiliaryMassModel}
  Idealization of the masses of the tether and balloon envelope
\end{definition}
\begin{proof}
  This is the declared model definition of Auxiliary Mass Model.
\end{proof}
\begin{definition}[Air Interaction Model]
  \label{def:physics:phyx-mini-0234:phyxminiproblems-problemphyxmini0234-airinteractionmodel}
  \lean{PhyXMiniProblems.ProblemPhyXMini0234.AirInteractionModel}
  Interaction of the ambient air with the moving balloon
\end{definition}
\begin{proof}
  This is the declared model definition of Air Interaction Model.
\end{proof}
\begin{definition}[Figure Side]
  \label{def:physics:phyx-mini-0234:phyxminiproblems-problemphyxmini0234-figureside}
  \lean{PhyXMiniProblems.ProblemPhyXMini0234.FigureSide}
  Left and right sides explicitly distinguished in the supplied image
\end{definition}
\begin{proof}
  This is the declared model definition of Figure Side.
\end{proof}
\begin{definition}[Figure Sphere]
  \label{def:physics:phyx-mini-0234:phyxminiproblems-problemphyxmini0234-figuresphere}
  \lean{PhyXMiniProblems.ProblemPhyXMini0234.FigureSphere}
  Two spheres visible in the supplied image
\end{definition}
\begin{proof}
  This is the declared model definition of Figure Sphere.
\end{proof}
\begin{definition}[Figure String Segment]
  \label{def:physics:phyx-mini-0234:phyxminiproblems-problemphyxmini0234-figurestringsegment}
  \lean{PhyXMiniProblems.ProblemPhyXMini0234.FigureStringSegment}
  Three string segments visible in the supplied image
\end{definition}
\begin{proof}
  This is the declared model definition of Figure String Segment.
\end{proof}
\begin{definition}[Supplied Figure Component]
  \label{def:physics:phyx-mini-0234:phyxminiproblems-problemphyxmini0234-suppliedfigurecomponent}
  \lean{PhyXMiniProblems.ProblemPhyXMini0234.SuppliedFigureComponent}
  \uses{def:physics:phyx-mini-0234:phyxminiproblems-problemphyxmini0234-figuresphere, def:physics:phyx-mini-0234:phyxminiproblems-problemphyxmini0234-figurestringsegment}
  Qualitative components of the supplied image
\end{definition}
\begin{proof}
  This is the declared model definition of Supplied Figure Component.
\end{proof}
\begin{definition}[Supplied Figure Readout]
  \label{def:physics:phyx-mini-0234:phyxminiproblems-problemphyxmini0234-suppliedfigurereadout}
  \lean{PhyXMiniProblems.ProblemPhyXMini0234.SuppliedFigureReadout}
  \uses{def:physics:phyx-mini-0234:phyxminiproblems-problemphyxmini0234-lengthquantity, def:physics:phyx-mini-0234:phyxminiproblems-problemphyxmini0234-figureside, def:physics:phyx-mini-0234:phyxminiproblems-problemphyxmini0234-suppliedfigurecomponent}
  Primary-image data. The two horizontal half-spans carry the label L, while the dashed vertical separation of the spheres carries the label y
\end{definition}
\begin{proof}
  This is the declared model definition of Supplied Figure Readout.
\end{proof}
\begin{definition}[Balloon Pendulum Setup]
  \label{def:physics:phyx-mini-0234:phyxminiproblems-problemphyxmini0234-balloonpendulumsetup}
  \lean{PhyXMiniProblems.ProblemPhyXMini0234.BalloonPendulumSetup}
  \uses{def:physics:phyx-mini-0234:phyxminiproblems-problemphyxmini0234-lengthquantity, def:physics:phyx-mini-0234:phyxminiproblems-problemphyxmini0234-timequantity, def:physics:phyx-mini-0234:phyxminiproblems-problemphyxmini0234-volumequantity, def:physics:phyx-mini-0234:phyxminiproblems-problemphyxmini0234-massquantity, def:physics:phyx-mini-0234:phyxminiproblems-problemphyxmini0234-massdensityquantity, def:physics:phyx-mini-0234:phyxminiproblems-problemphyxmini0234-accelerationquantity, def:physics:phyx-mini-0234:phyxminiproblems-problemphyxmini0234-forcequantity, def:physics:phyx-mini-0234:phyxminiproblems-problemphyxmini0234-angularfrequencyquantity, def:physics:phyx-mini-0234:phyxminiproblems-problemphyxmini0234-angularvelocityquantity, def:physics:phyx-mini-0234:phyxminiproblems-problemphyxmini0234-angularaccelerationquantity, def:physics:phyx-mini-0234:phyxminiproblems-problemphyxmini0234-balloongas, def:physics:phyx-mini-0234:phyxminiproblems-problemphyxmini0234-tetheranchor, def:physics:phyx-mini-0234:phyxminiproblems-problemphyxmini0234-pendulumorientation, def:physics:phyx-mini-0234:phyxminiproblems-problemphyxmini0234-auxiliarymassmodel, def:physics:phyx-mini-0234:phyxminiproblems-problemphyxmini0234-airinteractionmodel, def:physics:phyx-mini-0234:phyxminiproblems-problemphyxmini0234-suppliedfigurereadout}
  The physical state of the stated balloon experiment. The finite-amplitude motion and the first-order small-oscillation model have separate acceleration, frequency, and period data. In particular, no period field is defined to have the requested numerical value
\end{definition}
\begin{proof}
  This is the declared model definition of Balloon Pendulum Setup.
\end{proof}
\begin{definition}[Matches Verbal Balloon Scenario]
  \label{def:physics:phyx-mini-0234:phyxminiproblems-problemphyxmini0234-matchesverbalballoonscenario}
  \lean{PhyXMiniProblems.ProblemPhyXMini0234.MatchesVerbalBalloonScenario}
  \uses{def:physics:phyx-mini-0234:phyxminiproblems-problemphyxmini0234-balloonpendulumsetup}
  Categorical assumptions in the verbal statement: a helium balloon and light string are tethered to the ground, the light envelope is neglected, and air exerts buoyancy but no drag or other aerodynamic force
\end{definition}
\begin{proof}
  This is the declared model definition of Matches Verbal Balloon Scenario.
\end{proof}
\begin{definition}[Matches Supplied Figure]
  \label{def:physics:phyx-mini-0234:phyxminiproblems-problemphyxmini0234-matchessuppliedfigure}
  \lean{PhyXMiniProblems.ProblemPhyXMini0234.MatchesSuppliedFigure}
  \uses{def:physics:phyx-mini-0234:phyxminiproblems-problemphyxmini0234-lengthinmeters, def:physics:phyx-mini-0234:phyxminiproblems-problemphyxmini0234-balloonpendulumsetup}
  Readout of the supplied image. It deliberately records the visible two-wall, two-sphere, three-segment geometry without asserting that this inconsistent geometry governs the ground-tethered balloon dynamics
\end{definition}
\begin{proof}
  This is the declared model definition of Matches Supplied Figure.
\end{proof}
\begin{definition}[Matches Problem Data]
  \label{def:physics:phyx-mini-0234:phyxminiproblems-problemphyxmini0234-matchesproblemdata}
  \lean{PhyXMiniProblems.ProblemPhyXMini0234.MatchesProblemData}
  \uses{def:physics:phyx-mini-0234:phyxminiproblems-problemphyxmini0234-lengthinmeters, def:physics:phyx-mini-0234:phyxminiproblems-problemphyxmini0234-densityinkilogramspercubicmeter, def:physics:phyx-mini-0234:phyxminiproblems-problemphyxmini0234-angularvelocityinradianspersecond, def:physics:phyx-mini-0234:phyxminiproblems-problemphyxmini0234-balloonpendulumsetup}
  Numerical quantities printed in the problem: L = 3.00 m, helium density 0.179 kg/m\textasciicircum{}3, and air density 1.20 kg/m\textasciicircum{}3. The release conditions encode a nonzero displacement and release from rest. Rather than inventing a numeric cutoff for the source's qualitative word "slightly", smallness is represented below by the positive-amplitude limit and the little-o linearization
\end{definition}
\begin{proof}
  This is the declared model definition of Matches Problem Data.
\end{proof}
\begin{definition}[Uses Standard Gravity]
  \label{def:physics:phyx-mini-0234:phyxminiproblems-problemphyxmini0234-usesstandardgravity}
  \lean{PhyXMiniProblems.ProblemPhyXMini0234.UsesStandardGravity}
  \uses{def:physics:phyx-mini-0234:phyxminiproblems-problemphyxmini0234-accelerationinmeterspersecondsquared, def:physics:phyx-mini-0234:phyxminiproblems-problemphyxmini0234-balloonpendulumsetup}
  The conventional terrestrial calibration g = 9.8 m/s\textasciicircum{}2, needed to select a numerical answer choice. It is separated from the quantities printed in the problem statement
\end{definition}
\begin{proof}
  This is the declared model definition of Uses Standard Gravity.
\end{proof}
\begin{definition}[Has Physical Balloon Parameters]
  \label{def:physics:phyx-mini-0234:phyxminiproblems-problemphyxmini0234-hasphysicalballoonparameters}
  \lean{PhyXMiniProblems.ProblemPhyXMini0234.HasPhysicalBalloonParameters}
  \uses{def:physics:phyx-mini-0234:phyxminiproblems-problemphyxmini0234-lengthinmeters, def:physics:phyx-mini-0234:phyxminiproblems-problemphyxmini0234-timeinseconds, def:physics:phyx-mini-0234:phyxminiproblems-problemphyxmini0234-volumeincubicmeters, def:physics:phyx-mini-0234:phyxminiproblems-problemphyxmini0234-massinkilograms, def:physics:phyx-mini-0234:phyxminiproblems-problemphyxmini0234-densityinkilogramspercubicmeter, def:physics:phyx-mini-0234:phyxminiproblems-problemphyxmini0234-accelerationinmeterspersecondsquared, def:physics:phyx-mini-0234:phyxminiproblems-problemphyxmini0234-forceinnewtons, def:physics:phyx-mini-0234:phyxminiproblems-problemphyxmini0234-angularfrequencyinradianspersecond, def:physics:phyx-mini-0234:phyxminiproblems-problemphyxmini0234-balloonpendulumsetup}
  Positivity and nondegeneracy assumptions for the physical experiment
\end{definition}
\begin{proof}
  This is the declared model definition of Has Physical Balloon Parameters.
\end{proof}
\begin{definition}[Satisfies Balloon Pendulum Laws]
  \label{def:physics:phyx-mini-0234:phyxminiproblems-problemphyxmini0234-satisfiesballoonpendulumlaws}
  \lean{PhyXMiniProblems.ProblemPhyXMini0234.SatisfiesBalloonPendulumLaws}
  \uses{def:physics:phyx-mini-0234:phyxminiproblems-problemphyxmini0234-lengthinmeters, def:physics:phyx-mini-0234:phyxminiproblems-problemphyxmini0234-timeinseconds, def:physics:phyx-mini-0234:phyxminiproblems-problemphyxmini0234-volumeincubicmeters, def:physics:phyx-mini-0234:phyxminiproblems-problemphyxmini0234-massinkilograms, def:physics:phyx-mini-0234:phyxminiproblems-problemphyxmini0234-densityinkilogramspercubicmeter, def:physics:phyx-mini-0234:phyxminiproblems-problemphyxmini0234-accelerationinmeterspersecondsquared, def:physics:phyx-mini-0234:phyxminiproblems-problemphyxmini0234-forceinnewtons, def:physics:phyx-mini-0234:phyxminiproblems-problemphyxmini0234-angularaccelerationinradianspersecondsquared, def:physics:phyx-mini-0234:phyxminiproblems-problemphyxmini0234-balloonpendulumsetup}
  Exact governing laws for the nonlinear balloon pendulum. The first four fields are the light-envelope mass relation, equality of helium and displaced-air volumes, Archimedes' buoyant force, and weight. The net upward force is their difference. Crucially, the tangential Newton law retains sin theta; it does not replace the nonlinear restoring torque by an exact linear law. The amplitude-indexed period family records the physical meaning of the small-amplitude limit without assigning that limit a numerical value
\end{definition}
\begin{proof}
  This is the declared model definition of Satisfies Balloon Pendulum Laws.
\end{proof}
\begin{theorem}[sine sub id is Little O at zero]
  \label{thm:physics:phyx-mini-0234:phyxminiproblems-problemphyxmini0234-sine-sub-id-islittleo-at-zero}
  \lean{PhyXMiniProblems.ProblemPhyXMini0234.sine_sub_id_isLittleO_at_zero}
  The mathematical first-order contract behind sin theta ≈ theta. This is a local asymptotic statement at theta = 0, not a global equality. It can be proved from Real.hasDerivAt\_sin (or strengthened quantitatively using Real.sin\_bound)
\end{theorem}
\begin{proof}
  The conclusion follows from the governing relations and figure readouts named in the statement of sine sub id is Little O at zero.
\end{proof}
\begin{definition}[Satisfies First Order Small Oscillation Model]
  \label{def:physics:phyx-mini-0234:phyxminiproblems-problemphyxmini0234-satisfiesfirstordersmalloscillationmodel}
  \lean{PhyXMiniProblems.ProblemPhyXMini0234.SatisfiesFirstOrderSmallOscillationModel}
  \uses{def:physics:phyx-mini-0234:phyxminiproblems-problemphyxmini0234-lengthinmeters, def:physics:phyx-mini-0234:phyxminiproblems-problemphyxmini0234-timeinseconds, def:physics:phyx-mini-0234:phyxminiproblems-problemphyxmini0234-massinkilograms, def:physics:phyx-mini-0234:phyxminiproblems-problemphyxmini0234-forceinnewtons, def:physics:phyx-mini-0234:phyxminiproblems-problemphyxmini0234-angularfrequencyinradianspersecond, def:physics:phyx-mini-0234:phyxminiproblems-problemphyxmini0234-angularaccelerationinradianspersecondsquared, def:physics:phyx-mini-0234:phyxminiproblems-problemphyxmini0234-balloonpendulumsetup}
  Laws of the separate first-order small-oscillation model. The first two equations are exact statements about the *linearized surrogate*, not about the nonlinear acceleration above. The residual identity explicitly connects the surrogate to the exact sine dynamics. Combined with sine\_sub\_id\_isLittleO\_at\_zero, it says that the omitted restoring-torque term is lower order than the retained term as the angular amplitude tends to zero.
\end{definition}
\begin{proof}
  This is the declared model definition of Satisfies First Order Small Oscillation Model.
\end{proof}
\begin{lemma}[small Oscillation Angular Frequency sq]
  \label{lem:physics:phyx-mini-0234:phyxminiproblems-problemphyxmini0234-smalloscillationangularfrequency-sq}
  \lean{PhyXMiniProblems.ProblemPhyXMini0234.smallOscillationAngularFrequency_sq}
  \uses{def:physics:phyx-mini-0234:phyxminiproblems-problemphyxmini0234-lengthinmeters, def:physics:phyx-mini-0234:phyxminiproblems-problemphyxmini0234-densityinkilogramspercubicmeter, def:physics:phyx-mini-0234:phyxminiproblems-problemphyxmini0234-accelerationinmeterspersecondsquared, def:physics:phyx-mini-0234:phyxminiproblems-problemphyxmini0234-angularfrequencyinradianspersecond, def:physics:phyx-mini-0234:phyxminiproblems-problemphyxmini0234-balloonpendulumsetup, def:physics:phyx-mini-0234:phyxminiproblems-problemphyxmini0234-matchesproblemdata, def:physics:phyx-mini-0234:phyxminiproblems-problemphyxmini0234-hasphysicalballoonparameters, def:physics:phyx-mini-0234:phyxminiproblems-problemphyxmini0234-satisfiesballoonpendulumlaws, def:physics:phyx-mini-0234:phyxminiproblems-problemphyxmini0234-satisfiesfirstordersmalloscillationmodel}
  The physical laws imply the effective small-oscillation angular frequency. This relation is derived by comparing the two acceleration laws at the nonzero release displacement and eliminating mass, volume, buoyancy, and weight
\end{lemma}
\begin{proof}
  The conclusion follows from the governing relations and figure readouts named in the statement of small Oscillation Angular Frequency sq.
\end{proof}
\begin{lemma}[balloon Small Amplitude Limit Period formula]
  \label{lem:physics:phyx-mini-0234:phyxminiproblems-problemphyxmini0234-balloonsmallamplitudelimitperiod-formula}
  \lean{PhyXMiniProblems.ProblemPhyXMini0234.balloonSmallAmplitudeLimitPeriod_formula}
  \uses{def:physics:phyx-mini-0234:phyxminiproblems-problemphyxmini0234-lengthinmeters, def:physics:phyx-mini-0234:phyxminiproblems-problemphyxmini0234-timeinseconds, def:physics:phyx-mini-0234:phyxminiproblems-problemphyxmini0234-densityinkilogramspercubicmeter, def:physics:phyx-mini-0234:phyxminiproblems-problemphyxmini0234-accelerationinmeterspersecondsquared, def:physics:phyx-mini-0234:phyxminiproblems-problemphyxmini0234-balloonpendulumsetup, def:physics:phyx-mini-0234:phyxminiproblems-problemphyxmini0234-matchesverbalballoonscenario, def:physics:phyx-mini-0234:phyxminiproblems-problemphyxmini0234-matchessuppliedfigure, def:physics:phyx-mini-0234:phyxminiproblems-problemphyxmini0234-matchesproblemdata, def:physics:phyx-mini-0234:phyxminiproblems-problemphyxmini0234-hasphysicalballoonparameters, def:physics:phyx-mini-0234:phyxminiproblems-problemphyxmini0234-satisfiesballoonpendulumlaws, def:physics:phyx-mini-0234:phyxminiproblems-problemphyxmini0234-satisfiesfirstordersmalloscillationmodel}
  Period formula predicted by the first-order, small-amplitude model, T = 2 pi sqrt (L rho\_He / (g (rho\_air - rho\_He))). This is an equality for smallAmplitudeLimitPeriod, whose relation to the nonlinear finite-amplitude period is the limit law above. It is not asserted as an exact finite-amplitude period formula
\end{lemma}
\begin{proof}
  The conclusion follows from the governing relations and figure readouts named in the statement of balloon Small Amplitude Limit Period formula.
\end{proof}
\begin{definition}[Answer Choice]
  \label{def:physics:phyx-mini-0234:phyxminiproblems-problemphyxmini0234-answerchoice}
  \lean{PhyXMiniProblems.ProblemPhyXMini0234.AnswerChoice}
  Labels of the four period choices printed in the problem
\end{definition}
\begin{proof}
  This is the declared model definition of Answer Choice.
\end{proof}
\begin{definition}[answer Choice In Seconds]
  \label{def:physics:phyx-mini-0234:phyxminiproblems-problemphyxmini0234-answerchoiceinseconds}
  \lean{PhyXMiniProblems.ProblemPhyXMini0234.answerChoiceInSeconds}
  \uses{def:physics:phyx-mini-0234:phyxminiproblems-problemphyxmini0234-answerchoice}
  Displayed period, in seconds, associated with each answer label
\end{definition}
\begin{proof}
  This is the declared model definition of answer Choice In Seconds.
\end{proof}
\begin{definition}[Rounds To Nearest Hundredth]
  \label{def:physics:phyx-mini-0234:phyxminiproblems-problemphyxmini0234-roundstonearesthundredth}
  \lean{PhyXMiniProblems.ProblemPhyXMini0234.RoundsToNearestHundredth}
  A real value rounds to rounded at the nearest hundredth of a second
\end{definition}
\begin{proof}
  This is the declared model definition of Rounds To Nearest Hundredth.
\end{proof}
% --- Archon declaration topology end ---
% --- Archon physics formalization source end ---
